\begin{table*}[t]
\centering
\caption{Cheapest-feasible-degree selection under the pre-specified rule minimize $L(k)$ s.t.\ $\Delta$F1$(k)\le 0.03$. The selected $k^*$ is the lowest degree of the evaluated family in every tested seed and label, so the utility constraint never binds (Table~\ref{tab:t4_degree_grid}). This is not an adaptive-selection result. Proxy cohort (MIMIC-IV-ED-derived).}
\label{tab:s1_k_selection}
\begin{tabular}{lccc}
\toprule
$k$ & $\Delta$F1 & Latency (s) & Feasible \\
\midrule
3 & 0.0016 & 0.66 $\pm$ 0.02 & \checkmark $\leftarrow k^*$ \\
5 & 0.0019 & 0.91 $\pm$ 0.02 & \checkmark \\
7 & 0.0022 & 1.20 $\pm$ 0.02 & \checkmark \\
9 & 0.0026 & n/a$^\dagger$ & \checkmark \\
\bottomrule
\end{tabular}
\\[2pt]
{\footnotesize $\dagger$ the $\Delta$F1 of $k=9$ is simulation only; the original evaluator cannot execute $k=9$ (8 levels needed, 7 available), and real execution at $k=9$ uses the coefficient-folding evaluator (Table~\ref{tab:t5_real_vs_simulated}).}
\end{table*}
